\section{Dynamic ratio error and transient common-mode leakage}
\label{sec:dynamic}

The cancellation conditions above are exact identities between transfer functions, evaluated at every frequency. This section derives what they imply when the measured resistance and the excitation current vary during the observation window. The model remains the buffered front end of Section~\ref{sec:model-plane}: the sample resistance modulates the source amplitude, and the ladder itself is linear and time-invariant.

\subsection{Matched channels: an exact weighted-average identity}
\label{sec:dynamic-matched}

Let the two loaded branch transfers be equal after direct-current gain normalization, $G_P=G_N=G$ with $G(0)=1$ and impulse response $g(t)$, and form the ratio readout
\begin{equation}\label{eq:ratio-readout}
\hat R(t)=R_s\,\frac{(g\!\ast\! IR)(t)}{(g\!\ast\! I)(t)} ,
\end{equation}
where $I(t)$ is the excitation current and $R(t)$ the sample resistance. Because $R(t)$ is constant with respect to the convolution variable,
\begin{equation}\label{eq:weighted-average}
\hat R(t)-R(t)=\frac{\bigl(g\!\ast\! I\,[R(\cdot)-R(t)]\bigr)(t)}{(g\!\ast\! I)(t)} .
\end{equation}
Equation~\eqref{eq:weighted-average} is exact for every waveform of $I(t)$, including steps and square-wave ripple with fast edges: \emph{any} excitation-current fluctuation cancels from the ratio when the loaded transfers match. The remaining error is a weighted average of the past resistance values. If $|{\rm d}R/{\rm d}t|\le L_R$ and $|I|\le I_{\max}$, then
\begin{equation}\label{eq:lag-bound}
|\hat R(t)-R(t)|\;\le\;\frac{L_R\,I_{\max}\,M_1}{\bigl|(g\!\ast\! I)(t)\bigr|},
\qquad
M_1=\int_0^\infty \!\!\tau\,|g(\tau)|\,{\rm d}\tau .
\end{equation}
The bound uses the first absolute moment because the impulse response of an $N\ge2$ ladder is not nonnegative; residues alternate in sign although all poles are real and negative. For the ladder class, $M_1\le\sum_k |r_k|/|p_k|^2$ in terms of the poles and residues of $1/D_g$, which the coefficients of Section~\ref{sec:general-order} supply directly.

\subsection{Mismatched channels: the current-slew cross term}
\label{sec:dynamic-mismatch}

When the loaded transfers differ, $G_P\ne G_N$, the direct-current normalization removes the static ratio error, and the ratio error acquires a term driven by the current slew rate. Writing $\tau=D'(0)/D(0)$ for the direct-current group delay of the loaded transfer (in $x$ units, $\tau=(A_1+gB_1)/(1+gB_0)$, scaled by $R_*C_*$), the first-order expression is
\begin{equation}\label{eq:mismatch-dynamic}
\hat R-R \;=\; \text{lag term of Eq.~\eqref{eq:lag-bound}} \;+\; R\,\frac{\dot I}{I}\,\Delta\tau
\;+\;O\!\left(\ddot I/I,\ \Delta M_1\right),
\qquad \Delta\tau=\tau_P-\tau_N .
\end{equation}
The cross term $R\,(\dot I/I)\,\Delta\tau$ is absent from every static analysis: a pair that cancels common-mode perfectly at each frequency can still convert excitation-current dynamics into a ratio error through the group-delay mismatch of the two channels. For the second-order unit branch at $g=0.1$, a $+10\%$ deviation of the output capacitor gives $\Delta\tau\approx1.67~\mu{\rm s}$; at a relative current slew $\dot I/I=800~{\rm s}^{-1}$ and $R=100~{\rm k}\Omega$ this is $133~\Omega$ of dynamic ratio error.

\subsection{Equivalence of static cancellation and dynamic immunity}
\label{sec:welding}

The static conditions of Section~\ref{sec:general-order} and the dynamic property of Section~\ref{sec:dynamic-matched} are one statement. At the operating load $g$ the following are equivalent:
\begin{enumerate}
\item $D_P=D_N$ (equal coefficient vectors, the static cancellation condition);
\item $g_P=g_N$ (equal impulse responses);
\item the ratio readout~\eqref{eq:ratio-readout} rejects every excitation-current waveform, in the sense of Eq.~\eqref{eq:weighted-average}.
\end{enumerate}
Statements 1 and 2 are equal polynomials of equal degree. Statement 2 implies 3 by the derivation of Eq.~\eqref{eq:weighted-average}; conversely, driving both channels with an impulse-pair input isolates $g_P-g_N$, so 3 implies 2. The matching-family dimensions of Table~\ref{tab:load-dimensions} therefore count not only static common-mode cancellation but also the families inside which the dynamic ratio readout is exactly immune to excitation dynamics.

\subsection{Transient common-mode leakage after a load change}
\label{sec:cm-pulse}

For the load-shift family matched at $g_0$ and observed at $g$, Eq.~\eqref{eq:general-load-shift} gives $G_{DC}=-(g-g_0)\Delta B/(D_P D_N)$ with $\Delta B(0)=0$. The zero-frequency leakage vanishes, yet the response is band-limited, so a common-mode \emph{step} of amplitude $V_c$ produces a finite differential pulse
\begin{equation}\label{eq:cm-pulse}
y(t)=V_c\,\mathcal L^{-1}\!\left[\frac{G_{DC}(s)}{s}\right],
\qquad
\|y\|_2^2=V_c^2\sum_{k,j}\frac{r_k r_j}{-(p_k+p_j)} ,
\end{equation}
with the residue formula for the pulse energy and a peak bounded by $\|g_{DC}\|_1$; the slowest pole of $D_P D_N$ governs the recovery time. For the P/T branches of Table~\ref{tab:branches} matched at $g_0=0.1$ and observed at $g=0.2$, the closed-form step response peaks at $9.2\times10^{-3}\,V_c$ after $28~\mu{\rm s}$ and decays with a $19~\mu{\rm s}$ tail (all four poles real, between $-2.69$ and $-0.52$ in $x$ units).

\subsection{Transient verification}
\label{sec:dynamic-verif}

Four transient simulations checked these predictions with \(R_*=1~\mathrm{k\Omega}\), \(C_*=10~\mathrm{nF}\), and a \(10~\mathrm{k\Omega}\) operating load. With matched branches, a \(5~\mathrm{kHz}\) square-wave current ripple of \(\pm20\%\) with \(10~\mathrm{ns}\) edges and a resistance ramp from \(10~\mathrm{k}\Omega\) to \(120~\mathrm{k}\Omega\) in \(10~\mathrm{ms}\) gave a maximum ratio error of \(320.93~\Omega\), inside the bound of Eq.~\eqref{eq:lag-bound} with a \(1.08\times\) margin; the transient solver and an independent fourth-order reference solution agreed to six digits. With an output capacitor deviating by \(+10\%\) and a \(1~\mathrm{kHz}\) triangular current, the measured mean error of \(121.9~\Omega\) matched the cross-term prediction \(R(\dot I/I)\Delta\tau=133.3~\Omega\) within nine percent, and the joint scenario amplified the error to \(4.9\times\) the matched baseline. The common-mode step pulse of Section~\ref{sec:cm-pulse} agreed with the closed form to a relative difference of \(1.2\times10^{-7}\).
